\section{Formulación del plan del proyecto de investigación}
\label{sec:plan}

Esta sección estructura el proyecto como plan ejecutable y verificable. Se
declaran las etapas, actividades y entregables asociados a cada objetivo
específico, los resultados esperados y efectivamente obtenidos, los hitos
verificables, los criterios de seguimiento, y la evaluación de viabilidad
técnica y contextual.

\subsection{Estructura de etapas, actividades y entregables}

El plan se organiza en cuatro etapas secuenciales que corresponden
biunívocamente a los cuatro objetivos específicos y a las cuatro fases del
proceso ABP. La tabla~\ref{tab:plan} establece la correspondencia entre
objetivo, actividades, entregable e indicador de cumplimiento.

\begin{table}[H]
\centering
\caption{Plan del proyecto: correspondencia entre objetivos, actividades,
entregables e indicadores de cumplimiento.}
\label{tab:plan}
\footnotesize
\begin{tabular}{@{}p{1.5cm}p{3.5cm}p{3.6cm}p{4.6cm}@{}}
\toprule
\textbf{Objetivo} & \textbf{Actividades} & \textbf{Entregable} &
\textbf{Indicador de cumplimiento}\\
\midrule
OE1 &
Revisión de literatura indexada; extracción de rangos operativos de ruido;
parametrización de $\mathbf{Q}$ y $\mathbf{R}$ &
Matriz de revisión bibliográfica y parámetros de covarianza documentados &
Matrices declaradas con unidades explícitas y justificación física trazable
(Anexo~\ref{anexo:matriz}; sección~\ref{sec:unidad})\\
\addlinespace
OE2 &
Formulación del balance de masa; discretización; verificación de
observabilidad &
Modelo en espacio de estados en tiempo discreto &
Rango de la matriz de observabilidad igual a 2, verificado
computacionalmente (sección~\ref{sec:unidad})\\
\addlinespace
OE3 &
Codificación del núcleo predictor-corrector; modularización; generación de
telemetría degradada &
Módulo Python operativo con registro estructurado &
Ejecución sin error y salidas CSV idénticas byte a byte entre corridas
(sección~\ref{sec:calidad})\\
\addlinespace
OE4 &
Simulación de Monte Carlo; cálculo de métricas; prueba de hipótesis; análisis
de sensibilidad &
Informe de resultados, figuras y registros de trazabilidad &
Reducción del MSE estadísticamente significativa con $\alpha = 0{,}05$
(sección~\ref{sec:resultados})\\
\bottomrule
\end{tabular}
\end{table}

\subsection{Resultados esperados y resultados obtenidos}

El plan declaró como resultado esperado la obtención de una reducción
sustantiva y estadísticamente significativa del error cuadrático medio en la
reconstrucción del volumen, respecto de la lectura SCADA cruda como línea base.
La tabla~\ref{tab:esperado-obtenido} contrasta lo proyectado con lo
efectivamente alcanzado.

\begin{table}[H]
\centering
\caption{Contraste entre resultados esperados y resultados obtenidos.}
\label{tab:esperado-obtenido}
\footnotesize
\begin{tabular}{@{}p{4.6cm}p{4.2cm}p{4.6cm}@{}}
\toprule
\textbf{Resultado esperado} & \textbf{Resultado obtenido} & \textbf{Estado}\\
\midrule
Reducción sustantiva del MSE del volumen frente a la línea base &
Reducción del \SI{89,57}{\percent} (de \num{4,97} a \SI{0,52}{\metre\tothe{6}}) &
Cumplido\\
\addlinespace
Significancia estadística de la mejora con $\alpha = 0{,}05$ &
$t(49) = 135{,}24$; $p < 0{,}001$; IC 95\,\% $[1{,}488;\ 1{,}533]$~\si{\metre\cubed} &
Cumplido y reforzado con contraste no paramétrico\\
\addlinespace
Reconstrucción del estado latente no sensado &
Reducción del \SI{98,17}{\percent} en el RMSE de la tasa de llenado &
Cumplido; excede lo proyectado\\
\addlinespace
Costo computacional compatible con operación en tiempo real &
Ciclo recursivo sobre matrices $2\times2$; \num{1000} pasos en el orden de
milisegundos &
Cumplido\\
\addlinespace
Verificación de márgenes de estabilidad paramétrica &
Malla de desajuste $\mathbf{Q}/\mathbf{R}$ de $5\times5$ factores; 22 de 25
combinaciones superan a la línea base &
Cumplido; no estaba previsto en el plan inicial\\
\bottomrule
\end{tabular}
\end{table}

\subsection{Hitos verificables y trazabilidad con el proceso ABP}

El seguimiento del plan se ancla en los hitos formales de la asignatura, cada
uno de los cuales produjo evidencia documental evaluada. La
tabla~\ref{tab:hitos} registra esa trazabilidad.

\begin{table}[H]
\centering
\caption{Hitos verificables del proceso ABP y evidencia asociada.}
\label{tab:hitos}
\footnotesize
\begin{tabular}{@{}p{1.7cm}p{4.0cm}p{3.4cm}p{4.2cm}@{}}
\toprule
\textbf{Fase} & \textbf{Hito} & \textbf{Evidencia} & \textbf{Estado y ajuste
derivado}\\
\midrule
Fase 1 & Delimitación del problema y formulación de objetivos &
Sumativa 1 &
Aprobado; ajuste pendiente en plan del proyecto y citación, cerrado en esta
entrega\\
\addlinespace
Fase 2 & Marco teórico y enfoque metodológico &
Sumativa 2 &
Aprobado; ajuste en formato de referencias y trazabilidad con el artículo\\
\addlinespace
Fase 3 & Diseño metodológico y levantamiento de datos &
Formativas 3 y 4; Sumativa 3 &
Aprobado; ajustes en consideraciones éticas, unidades y supuesto de
normalidad, cerrados en esta entrega\\
\addlinespace
Fase 4 & Consolidación, resultados y comunicación &
Informe final, artículo IEEE y presentación &
En ejecución\\
\bottomrule
\end{tabular}
\end{table}

\subsection{Criterios de seguimiento del plan}

El seguimiento se sostiene en cuatro criterios operativos. En cuanto al
\textbf{cumplimiento de plazos}, cada etapa se ancla a la fecha de la evaluación
sumativa correspondiente, lo que impone un calendario externo verificable. En
cuanto al \textbf{alcance}, la sección~\ref{sec:problema} declara exclusiones
explícitas —interfaces gráficas, bases de datos persistentes, hardware y
acoplamiento a protocolos industriales— cuyo respeto puede auditarse
directamente sobre los entregables. En cuanto a la \textbf{calidad}, los
criterios de rigor de la sección~\ref{sec:calidad} definen umbrales
verificables, en particular la reproducibilidad exacta bajo semilla fija. Y en
cuanto a la \textbf{incorporación de retroalimentación}, cada observación
docente recibida se registra con su ajuste correspondiente y la sección donde se
materializa (Anexo~\ref{anexo:ajustes}).

\subsection{Viabilidad técnica y contextual}

La \textbf{viabilidad técnica} está acreditada por la ejecución efectiva del
proyecto. El módulo requiere únicamente un computador personal estándar y
bibliotecas científicas de código abierto; la complejidad computacional del
ciclo recursivo sobre matrices de dimensión $2\times2$ es despreciable frente a
la capacidad de cualquier equipo contemporáneo, lo que sostiene la afirmación de
compatibilidad con operación en tiempo real.

La \textbf{viabilidad contextual} se resolvió mediante una decisión metodológica
explícita. Los acuerdos de no divulgación vigentes en la industria minera
impiden el acceso a telemetría SCADA histórica de faenas operativas, lo que
haría inviable un diseño basado en datos reales dentro de los plazos y las
condiciones de acceso de un proyecto de titulación. La generación de datos
sintéticos paramétricos calibrados con literatura indexada
(sección~\ref{sec:tecnicas}) convierte esa restricción en un diseño ejecutable,
al costo —declarado en las secciones~\ref{sec:etica} y~\ref{sec:discusion}— de
limitar el nivel de madurez tecnológica alcanzable a laboratorio y simulación.

La \textbf{viabilidad de dedicación} se ajusta a una carga nominal de diez horas
semanales durante el semestre académico, restricción que motivó el blindaje
deliberado del alcance funcional en torno al motor algorítmico y el descarte de
módulos periféricos.

\subsection{Relación del plan con el artículo científico}

Cada etapa del plan alimenta una sección específica del manuscrito en formato
IEEE desarrollado en paralelo. La correspondencia detallada, incluyendo el
estado de transferencia de cada componente, se documenta en el
Anexo~\ref{anexo:correspondencia}. Esta relación no es declarativa: ambos
productos comparten el mismo conjunto de referencias verificadas, las mismas
figuras generadas por el código anexo y los mismos valores numéricos, de modo
que cualquier discrepancia entre ellos sería directamente detectable por un
evaluador.
